%==================================================================
% BAB 3 (PROPOSAL) — RANCANGAN PENELITIAN
%
% Berkas ini hanya dipakai oleh proposal.tex. Bab 1 dan Bab 2 dipakai
% bersama dengan naskah Tugas Akhir, sehingga "yang dijanjikan" pada
% proposal sama persis dengan yang kelak dibuktikan pada Tugas Akhir.
%
% Label yang harus ada di Bab 1 (rumusan masalah):
%   rq:rancang, rq:otomasi, rq:efektivitas
%==================================================================

\chapter{RANCANGAN PENELITIAN}

\section{Metodologi Penelitian}
Bab ini menjelaskan rancangan penelitian yang akan dilaksanakan untuk menjawab
rumusan masalah pada \cref{rq:rancang}, \cref{rq:otomasi}, dan
\cref{rq:efektivitas}. Penelitian akan dilakukan melalui pengembangan sistem
menggunakan model \textit{Rapid Application Development} (RAD) yang terdiri atas
tahap \textit{Requirement Planning}, \textit{RAD Design Workshop}, dan
\textit{Implementation}, dengan pembatasan hak akses berbasis
\textit{Role-Based Access Control} (RBAC). Sistem yang dihasilkan kemudian
dievaluasi melalui \textit{black-box testing}, \textit{User Acceptance Testing}
(UAT), dan \textit{System Usability Scale} (SUS) untuk memperoleh gambaran
kelayakan fungsional dan penerimaan pengguna. Alur penelitian secara keseluruhan
ditunjukkan pada \cref{fig:alurproposal}.

\begin{figure}
\centering
\includegraphics[width=\linewidth,height=8cm,keepaspectratio]{metode}
\caption{Langkah Penelitian}
\label{fig:alurproposal}
\end{figure}

\subsection{Studi Literatur}
Pada tahap ini peneliti akan menelaah penelitian terdahulu dan landasan teori
yang relevan dengan sistem informasi manajemen, model RAD, RBAC,
\textit{black-box testing}, UAT, dan SUS, serta penelitian sistem monitoring
magang pada \cite{Rosanti2025} dan \cite{Asehan2026}.
Studi literatur dilakukan untuk memastikan bahwa model pengembangan, instrumen
pengujian, dan teknik analisis yang dipilih sesuai dengan tujuan penelitian.

\subsection{\textit{Requirement Planning}}
Pada tahap ini peneliti akan merumuskan tujuan dan kebutuhan sistem bersama
calon pengguna, yaitu Koordinator Magang, Kepala Program Studi, Dosen
Pembimbing, Pembimbing Eksternal, dan Mahasiswa. Kebutuhan dikumpulkan melalui
wawancara dan pengamatan terhadap alur magang yang berjalan saat ini, mencakup
pengajuan dan validasi pendaftaran, pelaksanaan dan bimbingan, serta penilaian
akhir. Keluaran tahap ini adalah daftar kebutuhan fungsional dan non-fungsional
serta identifikasi hak akses setiap peran. Sesuai karakteristik RAD, tahap ini
dialokasikan singkat karena penyempurnaan kebutuhan akan berlanjut pada tahap
berikutnya (\cite{Parlika2023}).

\subsection{\textit{RAD Design Workshop}}
Pada tahap ini peneliti akan menyusun rancangan sistem bersama pengguna dalam
sesi-sesi perancangan. Rancangan mencakup \textit{use case diagram},
\textit{activity diagram}, rancangan basis data, rancangan antarmuka, dan
matriks peran terhadap izin akses (\textit{role--permission}) sebagai dasar
penerapan RBAC. Rancangan yang telah disusun ditinjau oleh pengguna, kemudian
direvisi sampai sesuai dengan kebutuhan.

\subsection{\textit{Implementation}}
Pada tahap ini rancangan diimplementasikan menjadi aplikasi web yang responsif
dan disusun per modul sesuai ruang lingkup fungsional, yaitu pendaftaran magang,
pelaksanaan dan bimbingan, serta penilaian akhir. Pengujian dilakukan secara
berulang selama implementasi, dan sistem dinyatakan siap digunakan apabila telah
dinilai layak oleh peneliti dan pengguna (\cite{Parlika2023}).

\subsection{Penyusunan Instrumen}
Instrumen penelitian terdiri atas tiga bagian. Pertama, lembar skenario
\textit{black-box testing} dengan teknik \textit{state transition} yang memuat
kolom transisi, asal, aksi, dan tujuan, serta kolom skenario, \textit{output},
dan kesimpulan (Mintarsih, 2023). Lembar ini mencakup skenario akses yang
diizinkan dan yang ditolak untuk setiap peran guna menguji penerapan RBAC.
Kedua, kuesioner UAT yang disusun pada empat variabel, yaitu fungsionalitas
sistem, kinerja sistem, pengalaman dan tampilan antarmuka, serta efisiensi dan
produktivitas (\cite{AliyahAliyah2025}). Ketiga, kuesioner SUS yang
terdiri atas 10 pernyataan (\cite{Fiandhika2024}). Kuesioner UAT dan SUS
menggunakan Skala Likert lima poin, dan masing-masing dilengkapi kolom saran
terbuka.

\subsection{Populasi dan Sampel}
Populasi penelitian adalah seluruh calon pengguna sistem pada Program Studi
Sistem Informasi Unika Atma Jaya, yaitu mahasiswa, Koordinator Magang, Dosen
Pembimbing, Pembimbing Eksternal, dan Kepala Program Studi. Sampel akan dipilih
secara purposif dengan kriteria terlibat atau akan terlibat langsung dalam alur
magang dan telah mencoba sistem. Setiap peran diwakili oleh sedikitnya satu
responden, dengan mahasiswa sebagai kelompok terbesar. Jumlah responden akan
ditetapkan bersama pembimbing dengan mempertimbangkan ketersediaan pengguna;
sebagai acuan, \cite{AliyahAliyah2025} melibatkan 16 responden yang
mewakili seluruh peran pengguna sistem.

\subsection{Pengumpulan Data}
Pengumpulan data dilakukan dalam dua bagian. Pada \textit{black-box testing},
peneliti menjalankan seluruh skenario pada lembar pengujian dan mencatat
\textit{output} yang dihasilkan sistem. Pada UAT dan SUS, responden akan
menjalankan sistem sesuai tugas yang relevan dengan perannya, kemudian mengisi
kuesioner daring melalui Google Forms. Setiap butir dinilai dengan Skala Likert
lima poin sebagaimana disajikan pada \cref{tab:likertproposal}.

\begin{table}[H]
    \centering
    \caption{Skala Likert yang akan digunakan}
    \label{tab:likertproposal}
    \small
    \begin{tabular}{|c|l|}
        \hline
        \textbf{Skor} & \textbf{Pernyataan} \\ \hline
        1 & Sangat Tidak Setuju (STS) \\ \hline
        2 & Tidak Setuju (TS)         \\ \hline
        3 & Netral (N)                \\ \hline
        4 & Setuju (S)                \\ \hline
        5 & Sangat Setuju (SS)        \\ \hline
    \end{tabular}
\end{table}

\subsection{Teknik Analisis Data}
Data yang terkumpul akan dianalisis melalui tiga tahap.

Pertama, hasil \textit{black-box testing} dicatat sebagai berhasil atau tidak
berhasil pada setiap transisi. Tingkat keberhasilan dihitung dengan
\cref{eq:keberhasilan}. Kesalahan yang ditemukan akan diperbaiki, lalu skenario
terkait diuji ulang sebelum UAT dilakukan.
\begin{equation}
    \label{eq:keberhasilan}
    \mathrm{Tingkat\ keberhasilan} =
    \frac{\mathrm{jumlah\ skenario\ berhasil}}{\mathrm{jumlah\ seluruh\ skenario}}
    \times 100\%
\end{equation}

Kedua, hasil UAT dianalisis secara deskriptif. Nilai rata-rata (\textit{mean})
setiap butir dihitung dengan membagi total bobot jawaban dengan jumlah
responden, kemudian dikonversi menjadi persentase terhadap bobot maksimum
sebagaimana pada \cref{eq:uat}.
\begin{equation}
    \label{eq:uat}
    \mathrm{Persentase} = \frac{\mathrm{nilai\ mean}}{\mathrm{bobot\ maksimum}}
    \times 100\%
\end{equation}
Persentase setiap butir dirata-ratakan per variabel, lalu diinterpretasikan
sebagai berikut: 0--20 persen sangat kurang baik, 21--40 persen kurang baik,
41--60 persen cukup baik, 61--80 persen baik, dan 81--100 persen sangat baik
(\cite{AliyahAliyah2025}).

Ketiga, skor SUS setiap responden dihitung dengan \cref{eq:sus}, dengan $R_i$
adalah skor jawaban responden pada pernyataan ke-$i$.
\begin{equation}
    \label{eq:sus}
    \mathrm{Skor\ SUS} = \left( \sum_{i \in \{1,3,5,7,9\}} (R_i - 1)
    + \sum_{i \in \{2,4,6,8,10\}} (5 - R_i) \right) \times 2{,}5
\end{equation}
Skor seluruh responden dirata-ratakan, dan hasilnya dibandingkan dengan nilai
ambang 68 yang lazim digunakan sebagai batas \textit{usability} rata-rata. Karena
skor SUS lebih mencerminkan persepsi pengguna (\cite{Fiandhika2024}),
hasilnya akan dibandingkan dengan temuan \textit{black-box testing}, skor UAT,
dan saran terbuka responden.

\subsection{Penulisan Laporan}
Tahap akhir adalah penyusunan laporan yang merangkum seluruh proses dan hasil
penelitian menjadi naskah Tugas Akhir yang utuh.

\section{Matriks Keterkaitan Penelitian}
Agar penelitian tetap terarah, setiap rumusan masalah dipetakan terhadap data dan
teknik analisis yang direncanakan, sebagaimana ditunjukkan pada
\cref{tab:linkageproposal}.

\begin{table}[H]
    \centering
    \caption{Rencana keterkaitan penelitian}
    \label{tab:linkageproposal}
    \small
    \begin{tabularx}{\linewidth}{|c|X|X|}
        \hline
        \textbf{RQ} & \textbf{Data yang direncanakan} & \textbf{Teknik analisis} \\ \hline
        \cref{rq:rancang}     & Kebutuhan fungsional dan non-fungsional, rancangan sistem (\textit{use case}, \textit{activity diagram}, basis data, antarmuka), dan hasil pengujian \textit{state transition} & Analisis kebutuhan, pemodelan sistem, dan \textit{black-box testing} \\ \hline
        \cref{rq:otomasi}     & Alur pengajuan, validasi, bimbingan, dan penilaian magang; hasil skenario akses per peran; skor UAT variabel fungsionalitas dan efisiensi & Pemodelan alur kerja, \textit{black-box testing}, dan analisis deskriptif UAT \\ \hline
        \cref{rq:efektivitas} & Skor UAT (4 variabel) dan skor SUS (10 pernyataan) dari responden & Persentase dan interpretasi skor UAT, perhitungan skor SUS \\ \hline
    \end{tabularx}
\end{table}

\section{Rencana dan Jadwal Penelitian}
Rencana pelaksanaan penelitian disusun selama enam bulan sebagaimana ditunjukkan
pada \cref{tab:jadwal}. Sel yang diarsir menandakan bulan pelaksanaan tiap
kegiatan.

\definecolor{jadwalisi}{gray}{0.25}
\begin{table}[H]
    \centering
    \caption{Rencana dan jadwal penelitian}
    \label{tab:jadwal}
    \small
    \begin{tabular}{|c|l|*{6}{>{\centering\arraybackslash}p{0.6cm}|}}
        \hline
        \multirow{2}{*}{\textbf{No.}} & \multirow{2}{*}{\textbf{Kegiatan}}
          & \multicolumn{6}{c|}{\textbf{Bulan}} \\ \cline{3-8}
         & & 1 & 2 & 3 & 4 & 5 & 6 \\ \hline
        1. & Studi literatur                 & \cellcolor{jadwalisi} &                       &                       &                       &                       &                       \\ \hline
        2. & \textit{Requirement Planning}   & \cellcolor{jadwalisi} & \cellcolor{jadwalisi} &                       &                       &                       &                       \\ \hline
        3. & \textit{RAD Design Workshop}    &                       & \cellcolor{jadwalisi} & \cellcolor{jadwalisi} &                       &                       &                       \\ \hline
        4. & \textit{Implementation}         &                       &                       & \cellcolor{jadwalisi} & \cellcolor{jadwalisi} &                       &                       \\ \hline
        5. & Pengujian (\textit{black-box}, UAT, SUS) &              &                       &                       & \cellcolor{jadwalisi} & \cellcolor{jadwalisi} &                       \\ \hline
        6. & Analisis data                   &                       &                       &                       &                       & \cellcolor{jadwalisi} &                       \\ \hline
        7. & Penulisan laporan               &                       &                       &                       &                       & \cellcolor{jadwalisi} & \cellcolor{jadwalisi} \\ \hline
        8. & Sidang                          &                       &                       &                       &                       &                       & \cellcolor{jadwalisi} \\ \hline
    \end{tabular}
\end{table}